% Public website CV, updated November 2026.
% Keep this version aligned with _data/cv.yml; use only cleared public details.
\documentclass[10pt,a4paper]{article}
\usepackage[T1]{fontenc}
\usepackage[utf8]{inputenc}
\usepackage{lmodern}
\usepackage[margin=1.4cm]{geometry}
\usepackage{microtype}
\usepackage{enumitem}
\usepackage{tabularx}
\usepackage{titlesec}
\usepackage[hidelinks]{hyperref}
\hypersetup{pdftitle={Yilin Liu -- Public Academic CV},pdfauthor={Yilin Liu}}
\pagestyle{empty}
\setlength{\parindent}{0pt}
\setlength{\parskip}{0pt}
\setlist[itemize]{leftmargin=1.3em,itemsep=1.5pt,topsep=3pt,parsep=0pt,partopsep=0pt}
\titleformat{\section}{\large\bfseries}{}{0pt}{}[\titlerule]
\titlespacing*{\section}{0pt}{8pt}{4pt}
\newcommand{\entry}[3]{%
  \begin{tabularx}{\textwidth}{@{}Xr@{}}
    \textbf{#1} & \textit{#2}
  \end{tabularx}\\[-1pt]
  \textit{#3}
}
\newcommand{\award}[2]{%
  \begin{tabularx}{\textwidth}{@{}p{0.09\textwidth}X@{}}
    #1 & #2
  \end{tabularx}\par
}
\begin{document}
\begin{center}
  {\LARGE\bfseries Yilin Liu}\\[3pt]
  \href{mailto:yliu859@connect.hkust-gz.edu.cn}{yliu859@connect.hkust-gz.edu.cn}
  \quad $\vert$ \quad
  \href{https://yilin-610-c.github.io/}{yilin-610-c.github.io}\\[2pt]
  {\small Public academic CV $\cdot$ November 2026}
\end{center}
\vspace{-6pt}

\section*{Research Interests}
Efficient LLM inference and serving; low-precision quantization; GPU-aware system--algorithm co-design; AI agents for systems development and evaluation.

\section*{Education}
\entry{The Hong Kong University of Science and Technology (Guangzhou)}{2023--2027}
{B.Sc. in Data Science and Big Data Technology}
\begin{itemize}
  \item September 2023--June 2027 (expected); final-year undergraduate. \textbf{Dean's List for three semesters.}
\end{itemize}
\entry{University of California, Berkeley}{Fall 2025}
{Visiting Student, Berkeley Global Access Program}
\begin{itemize}
  \item \textbf{GPA: 3.914/4.000}; studies in Computer Architecture and Systems Programming.
\end{itemize}

\section*{Research Experience}
\entry{Hardware-Aware KV-Cache Quantization for Efficient LLM Serving}{May 2026--Present}
{HKUST(GZ); Supervisor: Prof. Zeyi Wen}
\begin{itemize}
  \item Investigating quantization error, activation distributions, and accuracy--execution trade-offs in low-bit KV caches; developing and evaluating quantization algorithms, calibration approaches, and scale selection.
  \item Implementing numerical diagnostics and ablations; contributing literature analysis, figures, and writing in collaboration with teammates responsible for GPU kernels, serving integration, and system benchmarking.
\end{itemize}
\entry{Reliable On-Device AI Agents and Real-World Evaluation}{Jun. 2026--Present}
{MBZUAI UGRIP 2026; Supervisors: Prof. Chun Jason Xue and Prof. Youcheng Sun}
\begin{itemize}
  \item Exploring demonstration-to-skill workflows and token-efficient execution; designing and prototyping a cross-app personal knowledge base in Python and Android.
  \item Deploying and testing cross-session memory on real devices; debugging benchmark setup, Android automation, memory integration, and result verification to improve evaluation reliability.
  \item Continuing collaboration after the summer internship, contributing prototyping and testing within the team's broader benchmark effort.
\end{itemize}
\entry{AgenticMath: Multi-Agent Mathematical Data Synthesis}{May 2025--May 2026}
{HKUST(GZ); Supervisor: Prof. Jiaheng Wei}
\begin{itemize}
  \item Reviewed mathematical data synthesis methods and engineered a multi-agent pipeline for seed selection, question reformulation, answer refinement, and quality evaluation.
  \item Fine-tuned 3B--8B LLMs and conducted in-domain and out-of-domain evaluations. The collaborative study demonstrated competitive reasoning performance with 30--60K synthesized examples on the evaluated benchmarks.
\end{itemize}

\section*{Final-Year Project}
\entry{Benchmarking AI Agents for Quantization Kernel Development}{In progress}
{HKUST(GZ)}
\begin{itemize}
  \item Investigating evaluation methodology for coding agents on quantization-related GPU kernel tasks. Potential dimensions include correctness, numerical accuracy, performance, robustness, and reproducibility; scope is under development, with no release or results reported.
\end{itemize}

\section*{Publication}
Xianyang Liu, \textbf{Yilin Liu}, Shuai Wang, Hao Cheng, Andrew Estornell, Yuzhi Zhao, Jun Shu, and Jiaheng Wei.
\href{https://openreview.net/forum?id=wrOXg95VoF}{\textit{AgenticMath: Enhancing LLM Reasoning via Agentic-based Math Data Generation}}.
\textbf{Accepted at the DATA-FM Workshop at ICLR 2026}.
\href{https://arxiv.org/abs/2510.19361}{[arXiv preprint]}

\section*{Honors and Awards}
\award{2026}{\textbf{Best Project Award}, MBZUAI UGRIP 2026; one of three awarded projects among 22 (June).}
\award{2026}{\textbf{2025 INFH--Lizhi Innovative Talent Scholarship}; one of 10 undergraduates selected annually; awarded March 2026.}
\award{2026}{\textbf{Outstanding Award}, Ascend AI Innovation Competition, Operator Challenge, HKUST(GZ) Special Session (June).}
\award{2025}{Second Prize, Blue Bridge Cup National Software and IT Talent Competition, Guangdong Division.}
\award{2024}{\textbf{iGEM Gold Medal and Best New Basic Part Nomination.}}
\award{2024}{Second Prize, 3rd Electrical \& Electronic Engineering Innovation Competition.}
\end{document}
